%% 04-track-summaries.tex — One paragraph per track

\section{Research Track Summaries}

\paragraph{Track A --- Overview and Guides (5 papers).}
Track A contains the documents you are reading now, plus the portfolio guide,
the synthesis metapaper, the replication package, and the policy brief.
These papers are written for general audiences---judges, legislators, and journalists---
and serve as entry points to the more technical tracks.
Key finding: the entire portfolio can be summarized in five numbers
(see Section~\ref{sec:findings}).

\paragraph{Track B --- Algorithm (18 papers).}
Track B is the technical core of the portfolio.
It introduces recursive bisection (B.1), proves that edge-weighting
improves compactness by 22\% (B.2), compares recursive and n-way approaches
(B.3, B.5), studies adaptive parameter selection (B.4), and establishes
the $O(n^{1.07})$ runtime complexity (B.6).
Later B papers explore solution-space sensitivity (B.7), ratio-optimal
bisection (B.8), dual population-area constraints (B.9), subdivision
alignment (B.10), regional apportionment (B.11), proportional constraints
(B.12), nested multi-chamber design (B.13), minority-opportunity bisection (B.14),
cross-census stability (B.15), convergence parameters (B.16), parameter
sensitivity (B.17), and multi-reapportionment stability (B.18).
The track as a whole demonstrates that a single algorithm family---bisection---
handles every redistricting scenario from a two-district state to all 435 seats.

\paragraph{Track C --- Validation (9 papers).}
Track C asks whether the results hold up.
It analyzes sensitivity to the modifiable areal unit problem (C.1),
validates results across three census cycles (C.2), studies temporal
stability of district boundaries (C.3), tracks longitudinal demographic
trends (C.4), quantifies the 62\% reduction in partisan efficiency gap (C.5),
and investigates user interpretation, uncertainty quantification, competitive
elections, and real-world adoption case studies (C.6--C.9).
The headline result: algorithmic maps are not just theoretically fair---
they are empirically fairer than enacted maps across every metric tested.

\paragraph{Track D --- Legal and VRA Compliance (6 papers).}
Track D provides the legal foundation.
Paper D.1 establishes the 42\% minority-population threshold that determines
where proportional majority-minority representation is geographically feasible.
Subsequent papers compare n-way and recursive approaches under the Voting Rights
Act (D.2), examine compactness tradeoffs (D.3), analyze legal implementation
pathways (D.4), and apply Gingles bloc-voting methodology (D.5).
Track D as a whole demonstrates that algorithmic redistricting not only complies
with the VRA but in most states produces \emph{more} majority-minority districts
than enacted plans.

\paragraph{Track E --- Extensions (7 papers).}
Track E tests the algorithm beyond standard congressional redistricting:
multi-member districts (E.1), county-based representation (E.2), a
national 435-district single-run approach (E.3), partisan-similarity
clustering (E.4), party-proportional seat allocation (E.5), international
applications (E.6), and lessons learned across all extensions (E.7).
The track shows that recursive bisection is a general-purpose tool,
not a solution narrowly fitted to one jurisdiction's rules.

\paragraph{Track F --- State Legislatures (6 papers).}
Track F scales the algorithm to state legislative chambers,
which have different seat counts, different compactness requirements,
and in many states, bicameral nesting constraints (state house districts
must nest inside state senate districts).
Papers F.1--F.6 cover all-50-state single chambers, bicameral nesting,
high-seat-count chambers ($k > 100$), state-specific criteria, compactness
comparisons between legislative and congressional districts, and
VRA compliance at the state level.
The key finding is that nesting is achievable algorithmically with no
degradation in population balance or compactness.

\paragraph{Track G --- Ensemble Analysis (5 papers).}
Track G situates algorithmic maps within the universe of all possible maps.
By generating thousands of random redistricting plans and comparing them
to algorithmic plans, Track G provides the strongest possible evidence of
neutrality: the algorithm produces plans that are statistically
indistinguishable from random draws on compactness and partisan metrics,
with better convergence diagnostics than Markov chain methods (G.4).
Track G closes the portfolio by answering the hardest question a skeptic
can ask: ``How do we know the algorithm didn't just accidentally pick
a partisan outcome?''  The ensemble answer is: we don't rely on
the accident---we show the distribution.
